\chapter{Вероятностное пространство}

\section{Основные определения}

\begin{definition}
	Результаты опытов или наблюдений будем называть \emph{событиями}.
\end{definition}

\begin{notation}
	$ \omega $ "--- исход, элементарное событие.
\end{notation}

\begin{notation}
	$ \Omega $ "--- множество всех элементарных событий.
\end{notation}

\begin{definition}
	Если $ A \cap B = \O $, то говорим, что $ A $ и $ B $ \emph{несовместны}.
\end{definition}

\begin{definition}
	Если $ A \sub B $, то говорят, что появление события $ A $ \emph{влечёт} появление события $ B $.
\end{definition}

\begin{notation}
	$ \sigma $-алгебру событий будем обозначать $ \Im $.
\end{notation}

\begin{definition}
	Пусть $ P : \Im $ удовлетворяет следующим условиям:
	\begin{enumerate}
		\item $ P(A) \ge 0 \quad \forall A \in \Im $;
		\item $ P(\Omega) = 1 $;
		\item $ A_1, A_2, \dotsc $ "--- не более чем счётный набор событий,
			$ A_i \cap A_j = \O $, то
			$$ P \Bigl( \bigcup_{i = 1}^\infty A_i \Bigr) = \sum_{i = 1}^\infty P(A_i) $$
	\end{enumerate}

	Будем говорить, что $ P $ является \emph{вероятностной (счётно-аддитивной) мерой} или
	\emph{вероятностью}.
\end{definition}

\begin{definition}
	Говорят, что $ A_1, A_2, \dotsc $ "--- \emph{полная группа несовместных событий},
	если их объединение представляет собой достоверное событие.
\end{definition}

\section{Свойства вероятностей}

\begin{props}
\item $ \ol A $ "--- событие, дополнительное к событию $ A $.
	Тогда $ P(\ol A) = 1 - P(A) $.
\item $ 0 \le P(A) \le 1 $.
\item $ P(\O) = 0 $.
\item $ P(A \cup B) = P(A) + P(B) - P(AB) $
	$$ P \Bigl( \bigcup_{i = 1}^n A_i \Bigr) =
	\sum_{i = 1}^n P(A_i) - \sum_{i < j} P(A_iA_j) + \sum_{i < j < k}P(A_iA_jA_k) -
	\dots + (-1)^{n + 1}P \Bigl( \bigcap_{i = 1}^n A_i \Bigr) $$
\item $ P(\bigcup A_i) \le \sum P(A_i) $.
\item $ A \sub B \implies P(A) \le P(B) $.
\item $ A_1 \sub A_2 \sub \dotso, \quad A = \bigcup A_n $
	$$ P(A) = \lim\limits_{n \to \infty} P(A_n) $$
\item $ B_1 \supset B_2 \supset \dotso, \quad B = \bigcap B_n $
	$$ P(B) = \lim\limits_{n \to \infty} P(B_n) $$
\item $ A_1, A_2, \dotsc $ "--- не более чем счётная полная группа событий
	$$ \sum P(A_n) = 1 $$
\end{props}

\begin{statement}[неравенства Бонферрони]
	$$ \sum_{k = 1}^n P(A_k) - \sum_{1 \le k \le l \le n} P(A_kA_l) \le
	P \Bigl( \bigcup_{k = 1}^n A_k \Bigr) \le \sum_{k = 1}^n P(A_k) $$
\end{statement}

\section{Условные вероятности}

\begin{definition}
	$ A, B \in \Im, \quad P(B) > 0 $

	\emph{Условной вероятностью} $ P(A|B) $ называется отношение
	$$ P(A|B) = \frac{P(AB)}{P(B)} $$
\end{definition}

\begin{statement}
	Условная вероятность является вероятностью.
\end{statement}

\begin{eproof}
\item
	$$ P(A|B) = \frac{P(AB)}{P(B)} \ge 0 $$
\item
	$$ P(\Omega|B) = \frac{P(\Omega B)}{P(B)} = \frac{P(B)}{P(B)} = 1 $$
\item
	$$ P \Bigl( \bigcup_k A_k \Bigm| B \Bigr) =
	\frac{P \Bigl\lgroup B \cap \Bigl( \bigcup_k A_k \Bigr) \Bigr\rgroup}{P(B)} =
	\frac{P \Bigl( \bigcup_k A_kB \Bigr)}{P(B)} =
	\frac{\sum_k P(A_kB)}{P(B)} = \sum_k P(A_k|B) $$
\end{eproof}

\section{Формула полной вероятности}

\begin{theorem}
	$ A_1, A_2, \dotsc $ "--- полная группа несовместных событий, $ \quad P(A_i) > 0 $

	$$ P(B) = \sum_i P(B|A_i) \cdot P(A_i) $$
\end{theorem}

\begin{proof}
	Учитывая несовместность событий $ BA_1, BA_2, \dotsc $, получаем
	\begin{multline*}
		P(B) = P(B \cap \Omega) =
		p \Bigl\lgroup B \cap \Bigl( \bigcup_{i \ge 1}A_i \Bigr) \Bigr\rgroup =
		P \Bigl( \bigcup_{i \ge 1} BA_i \Bigr) = \sum_{i \ge 1} P(BA_i) = \\
		= \sum_{i \ge 1} \frac{P(BA_i)P(A_i)}{P(A_i)} = \sum_{i \ge 1} P(B|A_i) P(A_i)
	\end{multline*}
\end{proof}

\begin{theorem}[формула Байеса]
	$ A_1, A_2, \dotsc $ "--- полная группа несовместных событий \\
	$ P(A_i) > 0, \quad P(B) > 0 $

	$$ P(A_i|B) = \frac{P(B|A_i) \cdot P(A_i)}{\sum_j P(B|A_j) \cdot P(A_j)} $$
\end{theorem}

\begin{proof}
	$$ \frac{P(B|A_i)P(A_i)}{\sum_j P(B|A_j)P(A_j)} = \frac{P(B|A_i)P(A_i)}{P(B)} =
	\frac{P(A_iB)P(A_i)}{P(A_i)P(B)} = P(A_i|B) $$
\end{proof}

\section{Независимость случайных событий}

\begin{definition}
	$ P(B) \ne 0 $

	События $ A $ и $ B $ будем называть \emph{независимыми}, если $ P(A|B) = P(A) $.
\end{definition}

\begin{definition}
	События $ A $ и $ B $ будем называть \emph{независимыми}, если при $ P(A)P(B) \ne 0 $ выполнены
	равенства
	$$ P(A|B) = P(A), \quad P(B|A) = P(B) $$
\end{definition}

\begin{definition}
	События $ A $ и $ B $ называются \emph{независимыми}, если $ P(AB) = P(A)P(B) $.
\end{definition}

\begin{definition}
	События $ A_1, A_2, \dotsc $ \emph{попарно независимы}, если для любых двух событий $ A_i $ и
	$ A_j $ выполняется соотношение
	$$ P(A_iA_j) = P(A_i)P(A_j) $$
\end{definition}

\begin{definition}
	Говорят, что события $ A_1, A_2, \dots, A_n $ \emph{независимы в совокупности (взаимно
	независимы)}, если
	$$ \forall k \le n \quad \forall 1 \le i_1 < i_2 < \dots < i_k \le n \quad
	P \Bigl( \bigcap_{j = 1}^k A_{i_j} \Bigr) = \prod_{j = 1}^k P(A_{i_j}) $$
\end{definition}

\begin{theorem}
	$ A_1, A_2, \dotsc $ взаимно независимы.

	Тогда взаимно независимы события $ B_1, B_2, \dotsc $, где $ B_i $ "--- это $ A_i $ или его
	дополнение.
\end{theorem}

\section{Схемы испытаний Бернулли}

\begin{definition}
	Повторные независимые испытания называются \emph{испытаниями Бернулли}, если каждое такое
	испытание имеет только два возможных исхода и вероятности соответствующих исходов остаются
	неизменными для всех испытаний.
\end{definition}

\begin{statement}
	Вероятность из $ n $ испытаний получить ровно $ m $ успехов равна
	$$ P_n(m) = C_n^m p^mq^{n - m} $$
\end{statement}

\begin{theorem}[локальная предельная]
	$ 0 < p < 1 $

	Соотношение
	$$ \frac{\sqrt{npq}P_n(m)}{\phi \Bigl( \dfrac{m - np}{\sqrt{npq}} \Bigr)}
	\underarr{n \to \infty} 1, $$
	где
	$$ \phi(x) = \frac1{\sqrt{2\pi}} e^{-\frac{x^2}2}
	\text{ "--- плотность нормального распределения}, $$
	выполняется равномерно для значений $ m $ таких, что
	$$ -\infty < C_1 \le \frac{m - np}{\sqrt{npq}} \le C_2 < +\infty $$
\end{theorem}

\begin{notation}
	$ \mu_n $ "--- количество успехов в испытаниях Бернулли.
\end{notation}

\begin{notation}
	Определим \emph{функцию распределения нормального закона}:
	$$ \Phi(x) = \frac1{\sqrt{2\pi}} \int_{-\infty}^x e^{-\frac{t^2}2} \di t =
	\int_{-\infty}^x \phi(t) \di t $$
\end{notation}

\begin{theorem}[интегральная предельная]
	$ 0 < p < 1 $

	Равномерно по всем значениям $ -\infty \le a < b \le +\infty $ выполняется соотношение
	$$ P\set{a \le \mu_n \le b} \underarr{n \to \infty}
	\Phi \Bigl( \frac{b - np}{\sqrt{npq}} \Bigr) - \Phi \Bigl( \frac{a - np}{\sqrt{npq}} \Bigr) $$
\end{theorem}

Погрешность оценивается неравенством
$$ \Delta_n \le \frac c{\sqrt{npq}} $$
$ c $ "--- абсолютная константа.
Можно взять $ c = 1 $ или $ c = 0.7655 $.

\begin{theorem}
	При большом $ n $ рассмотрим такую схему серий независимых событий: \\
	\begin{tabular}{l | c}
		& p \\
		\hline
		$ A_{11} $ & \lambda \\
		$ A_{21} $ $ A_{22} $ & \lambda/2 \\
		$ A_{31} $ $ A_{32} $ $ A_{33} $ & \lambda/3 \\
		\vdots & \vdots \\
		$ A_{n1} $ $ A_{n2} $ \dots $ A_{nn} $ & \lambda/n
	\end{tabular}

	Для каждого фиксированного $ m $ справедливо
	$$ P_n(m) \underarr{n \to \infty} \frac{e^{-\lambda}\lambda^m}{m!} $$
\end{theorem}

\section{Полиномиальная схема}

\begin{definition}
	Пусть $ P\Set{A_i} = p_i, \quad p_1 + \dots + p_r = 1 $.

	Вероятность того, что в $ n $ испытаниях событие $ A_i $ произойдёт $ m_i $ раз, равна:
	$$ P_n(m_1, m_2, \dots, m_r) =
	\frac{n!}{m_! m_2! \cdots m_r!} p_1^{m_1} p_2^{m_2} \cdots p_r^{m_r}, \quad
	m_1 + m_2 + \dots + m_r = n $$
\end{definition}

